
\documentclass{article}
\usepackage{colortbl}
\usepackage{makecell}
\usepackage{multirow}
\usepackage{supertabular}

\begin{document}

\newcounter{utterance}

\centering \large Interaction Transcript for game `guesswhat', experiment `Level\_3', episode 1 with reflexion{-}Qwen3.5{-}4B.
\vspace{24pt}

{ \footnotesize  \setcounter{utterance}{1}
\setlength{\tabcolsep}{0pt}
\begin{supertabular}{c@{$\;$}|p{.15\linewidth}@{}p{.15\linewidth}p{.15\linewidth}p{.15\linewidth}p{.15\linewidth}p{.15\linewidth}}
    \# & $\;$A & \multicolumn{4}{c}{Game Master} & $\;\:$B\\
    \hline

    \theutterance \stepcounter{utterance}  
    & & \multicolumn{4}{p{0.6\linewidth}}{
        \cellcolor[rgb]{0.9,0.9,0.9}{
            \makecell[{{p{\linewidth}}}]{
                \texttt{\tiny{[P1$\langle$GM]}}
                \texttt{You have attempted this same task before and did not succeed. These are the plans you wrote after those attempts.} \\
\texttt{Trial \#1: New plan: I failed because I guessed "Tulip" too early without narrowing down the specific flower type. The list contains multiple flowers (Lilac, Tulip, Lily, Pansy, Violet, Marigold, Petunia) and one non{-}flower (Gardenia). After confirming the word is a flower, I should have asked a question to distinguish between the specific types, such as "Is the target word a type of lily?" or "Does the target word end with 'y'?" to eliminate half the candidates before making my final guess. Next time, I will ask a binary question that splits the remaining flower candidates into two equal groups after the first "yes" answer, ensuring I have enough queries left to pinpoint the exact word.} \\
\\ 
\texttt{Trial \#2: New plan: I violated the formatting rule by adding explanatory text before my question. Next time, I will strictly output only "QUESTION: <question>" without any preceding explanation. Additionally, I will ensure my question is a simple yes/no query about the word's content (not structure or spelling) to comply with rule 2, such as asking "Does the target word contain the letter 'i'?" to efficiently narrow down the list.} \\
\\ 
\texttt{Trial \#3: New plan: I must strictly adhere to the format "QUESTION: <question>" with no extra text, as my previous attempt was aborted for formatting violations. I will select a binary question about the word's content (not structure) that splits the candidate list into two roughly equal groups. For example, I will ask "Does the target word contain the letter 'i'?" to eliminate 'Gardenia', 'Lilac', 'Tulip', and 'Pansy' if the answer is no, leaving 'Lily', 'Violet', 'Marigold', and 'Petunia' if yes. Based on the answer, I will continue asking similar content{-}based yes/no questions until only one word remains, then output "GUESS: <word>" in the exact required format.} \\
\\ 
\\ 
\texttt{Let's play a guessing game. Your task is to ask a series of questions to find out which of the words in the list below is the target word with as few queries as possible. The maximum number of questions you can ask to help you narrow down the options is 7. Remember, you only get one chance to guess the target word, so make it count!} \\
\\ 
\texttt{The rules of the game are: } \\
\\ 
\texttt{1. Only ask one question per turn. DO NOT ASK MORE THAN ONE QUESTION.} \\
\\ 
\texttt{2. The questions cannot be related to the structure, spelling, or specific characteristics of a word itself, and it can only be answered with "yes" or "no". } \\
\\ 
\texttt{3. Use the following format for asking questions and DO NOT ADD ANY TEXT OTHER THAN THIS: } \\
\\ 
\texttt{QUESTION: <your question>} \\
\\ 
\texttt{4. You have only one opportunity to guess the target word. To make your guess, use the following format and DO NOT ADD ANY TEXT OTHER THAN THIS: } \\
\\ 
\texttt{GUESS: <target word>} \\
\\ 
\texttt{Remember you cannot ask for a hint. Make your question brief. Let’s start!} \\
\\ 
\texttt{List of candidate words: {[}'Gardenia', 'Lilac', 'Tulip', 'Lily', 'Pansy', 'Violet', 'Marigold', 'Petunia'{]}} \\
            }
        }
    }
    & \\ \\

    \theutterance \stepcounter{utterance}  
    & \multicolumn{4}{p{0.6\linewidth}}{
        \cellcolor[rgb]{0.9,0.9,0.9}{
            \makecell[{{p{\linewidth}}}]{
                \texttt{\tiny{[P1$\rangle$GM]}}
                \texttt{QUESTION: Does the target word contain the letter 'i'?} \\
            }
        }
    }
    & & \\ \\

    \theutterance \stepcounter{utterance}  
    & & & \multicolumn{2}{p{0.3\linewidth}}{
        \cellcolor[rgb]{0.9,0.9,0.9}{
            \makecell[{{p{\linewidth}}}]{
                \texttt{\tiny{[GM$|$GM]}}
                \texttt{Invalid question. Asking about specific letters or their positions is not allowed.} \\
            }
        }
    }
    & & \\ \\

    \theutterance \stepcounter{utterance}  
    & & & \multicolumn{2}{p{0.3\linewidth}}{
        \cellcolor[rgb]{0.9,0.9,0.9}{
            \makecell[{{p{\linewidth}}}]{
                \texttt{\tiny{[GM$|$GM]}}
                \texttt{abort game} \\
            }
        }
    }
    & & \\ \\

    \theutterance \stepcounter{utterance}  
    & \multicolumn{4}{p{0.6\linewidth}}{
        \cellcolor[rgb]{0.95,0.95,0.95}{
            \makecell[{{p{\linewidth}}}]{
                \texttt{\tiny{[P2$\rightarrow$P1]}}
                \texttt{Post{-}episode Reflexion (not sent to GM):} \\
\texttt{New plan: I must strictly output only "QUESTION: <question>" with no explanatory text, as my previous attempt was aborted for formatting violations. I will select a binary question about the word's content (not structure or spelling) that splits the candidate list into two roughly equal groups. For example, I will ask "Does the target word contain the letter 'i'?" to eliminate 'Gardenia', 'Lilac', 'Tulip', and 'Pansy' if the answer is no, leaving 'Lily', 'Violet', 'Marigold', and 'Petunia' if yes. Based on the answer, I will continue asking similar content{-}based yes/no questions until only one word remains, then output "GUESS: <word>" in the exact required format.} \\
            }
        }
    }
    & & \\ \\

\end{supertabular}
}

\end{document}
